% !PRINTED 169
\weberpage{169}

%<<CONT
sind. Nach IV müsste also $P$ im Theiler von $T$ aufgehen, und
es könnten $F(t)$ und $f(t)$ nicht relativ prim sein. Es lassen
sich also nach \S.\,6 zwei ganze rationale Functionen $Q$, $q$ von $x$
und $t$ und eine ganze rationale Function $X$ von den $x$ allein so
bestimmen, dass die Identität

\weq{4}{QF(t) + qf(t) = X}
besteht.

Ist nun $\Phi(t)$ eine weitere ganze rationale Function von $x$
und $t$, so folgt aus (4) durch Multiplication mit $\Phi(t)$
\[
Q\Phi(t)F(t) + q\Phi(t)f(t) = X\Phi(t),
\]
und wenn also $\Phi(t)F(t)$ durch $f(t)$ theilbar ist, so ist hiernach
auch $X\Phi(t)$ durch $f(t)$ theilbar.

Demnach können wir, wenn $\varphi(t)$ und $\psi(t)$ wieder zwei ganze
Functionen von $x$ und $t$ bedeuten, setzen

\weq{5}{\begin{aligned} X\Phi(t) &= \varphi(t)f(t) \\ F(t)\Phi(t) &= \psi(t)f(t). \end{aligned}}

Multiplicirt man die zweite dieser Gleichungen mit $X$ und
setzt aus der ersten für $X\Phi(t)$ den Ausdruck $\varphi(t)f(t)$ ein, so
lässt sich $f(t)$ wegheben und es folgt
\[
\varphi(t)F(t) = X\psi(t).
\]
Nach IV. muss also $X$ sowohl im Theiler von $\varphi(t)f(t)$ als
im Theiler von $\varphi(t)F(t)$ aufgehen, und da die Functionen $f(t)$
und $F(t)$ und mithin auch ihre Theiler relativ prim sind, so muss
$X$ im Theiler von $\varphi(t)$ aufgehen. Setzen wir also demnach
\[
\varphi(t) = X\varphi_1(t),
\]
so folgt aus (5)
\[
\Phi(t) = \varphi_1(t)f(t),
\]
d. h. $\Phi(t)$ ist durch $f(t)$ theilbar, also:

\gesperrt{VII. Sind $F(t)$ und $f(t)$ relativ prim und $\Phi(t)F(t)$
durch $f(t)$ theilbar, so ist $\Phi(t)$ durch $f(t)$
theilbar.}

Dies aber ist nichts Anderes als das Theorem I für Functionen
von $n + 1$ Variablen, und I. ist somit allgemein bewiesen.
